\chapter{PEM Fuel Cell Solver}
\label{s:pemfc}

\section{Synopsis}

The PEMFCSolver simulates species transport and electrochemistry in
proton exchange membrane (PEM) fuel cells on two-dimensional
channel-plane domains. Channel and rib layouts are represented on a
fixed mesh through Brinkman penalization, so no remeshing is needed
when the layout changes; this also makes the solver directly usable
in density-based topology optimization workflows. The flow field is
computed separately by the IncNavierStokesSolver, using the
\inltt{VCSImplicitAbsorption} solver type described below, and is
passed to the species solver as a field file.

The solver family covers the following equation types.

\begin{table}[h!]
\begin{center}
\small
\renewcommand\arraystretch{1.6}
\begin{tabular}{lp{8.5cm}}
\toprule
\textbf{EquationType} & \textbf{Description} \\
\midrule
\inltt{PEMFCSpecies} & Unsteady advection-diffusion of oxygen with a
nonlinear Butler-Volmer sink, IMEX time integration; the local
current density is resolved pointwise by a safeguarded Newton
iteration and written to output as a derived field. \\
\inltt{PEMFCSteady} & Steady counterpart: Picard (secant) iteration
in which the Butler-Volmer sink is linearized about the previous
iterate; each iterate is one variable-coefficient steady
advection-diffusion-reaction solve. \\
\inltt{PEMFCCathodeMS} & Steady cathode multi-species system
(oxygen and water vapor, nitrogen closing the mixture) with local
mixture diffusivities, a derived water source ratio and an optional
Stefan correction. \\
\inltt{PEMFCMultilayer} & Steady multilayer (section-slice) system:
several stacked physical layers collapsed onto one 2D mesh, coupled
by volumetric inter-layer exchange terms, with a pointwise
two-electrode kinetics closure. \\
\bottomrule
\end{tabular}
\end{center}
\end{table}

In addition, the incompressible Navier-Stokes solver gains the
solver type \inltt{VCSImplicitAbsorption}: a velocity correction
scheme in which a spatially varying absorption (Brinkman) term
$-\sigma(\mathbf{x})\,\mathbf{u}$ is treated implicitly inside the
velocity Helmholtz solves. This removes the explicit-forcing
stability restriction $\sigma \Delta t \lesssim 1$, which is severe
at realistic penalization strengths ($\sigma \sim 10^6$~s$^{-1}$).

\section{Usage}

\begin{lstlisting}[style=BashInputStyle]
PEMFCSolver session.xml
\end{lstlisting}

The flow field is computed first, for example
\begin{lstlisting}[style=BashInputStyle]
IncNavierStokesSolver mesh.xml flow.xml
\end{lstlisting}
with \inltt{SolverType} set to \inltt{VCSImplicitAbsorption} and the
absorption coefficient given by the \inltt{ImplicitAbsorption}
function. The converged velocity is then supplied to the species
session through the \inltt{AdvectionVelocity} function with a
\inltt{FILE} attribute.

\section{Session file configuration}

\subsection{Solver info}

\begin{itemize}
\item \inltt{EqType}: one of the equation types listed above.
\item \inltt{Projection}: \inltt{Continuous} (the solver family is
formulated for continuous Galerkin projections).
\end{itemize}

\subsection{Parameters}

The main physical parameters of the steady cathode solvers:

\begin{itemize}
\item \inltt{epsilon}: effective diffusivity of the transported
species.
\item \inltt{Htot}: total out-of-plane height entering the
volumetric sink.
\item \inltt{Vcell}: cell voltage.
\item \inltt{Rlump}: lumped area-specific resistance of the
Butler-Volmer closure.
\item \inltt{i0}: exchange current density.
\item \inltt{WaterSourceRatio}: mass ratio of the net water source
to the oxygen sink (\inltt{PEMFCCathodeMS}).
\item \inltt{StefanCorrection}: set to 1 to enable the Stefan flow
correction of the diffusive fluxes (\inltt{PEMFCCathodeMS}).
\item \inltt{PicardIters}, \inltt{PicardTol}, \inltt{PicardRelax}:
outer iteration controls of the steady solvers.
\end{itemize}

\subsection{Functions}

\begin{itemize}
\item \inltt{AdvectionVelocity}: the advecting velocity, normally
read from a field file produced by the flow solve.
\item \inltt{DiffusivityFactor}: spatially varying factor applied to
the diffusivity, used to distinguish channel and rib regions.
\item \inltt{ImplicitAbsorption} (flow sessions): the absorption
coefficient $\sigma(\mathbf{x})$ of
\inltt{VCSImplicitAbsorption}.
\end{itemize}

Note for sessions written for Nektar++ 5.0: the flow sessions must
use the \inltt{TIMEINTEGRATIONSCHEME} block (for example method
\inltt{IMEX} of order 2) instead of the removed
\inltt{TimeIntegrationMethod} solver-info property.

\section{Examples}

The regression test in \inlsh{solvers/PEMFCSolver/Tests} solves a
cathode two-species configuration on a coarse mesh of the parallel
flow field verification case of the solver publication: flow through
a penalized channel-rib layout, oxygen and water-vapor transport,
and Butler-Volmer kinetics with a Picard steady iteration. The
converged fields reproduce the published verification values.
